\chapter{Design}\label{s:design}

This chapter motivates the design decisions behind the architecture presented
in Chapter~\ref{s:overview}. Each decision is presented together with the
alternative it was chosen over and the evidence, where available, that
justified the choice. 

\section{Intent Decomposition}\label{s:decomposition}

Before any retrieval or generation takes place, the pipeline runs the raw
intent through a dedicated decomposition step: a short,
separate prompt asks the same already-loaded generation model to break the
intent into a numbered list of at most <N> concrete sub-tasks, phrased in
plain English rather than code. 

This step was introduced because retrieving directly against the raw intent
tried to answer two different questions at once:
understanding \emph{what needs to be done} (which BraneScript constructs and
package calls the request actually requires) and \emph{how to express it}
(what BraneScript syntax to use). A single embedding of a multi-part
intent tends to average over all of its sub-goals rather than surfacing a
chunk for each one. Stating explicitly, in English, that a request needs a
package call, a loop over a dataset, and execution on a named node gives the
retrieval stage (Section~\ref{s:retrieval-design}) several focused queries
instead of one broad one; the per-subtask fallback search described there
only exists because this decomposition happens first, sharpening the query
before it is ever embedded, thereby directly shaping the
retrieval-augmented context whose quality is examined in~\textbf{RQ2}.

Two constraints shape what a sub-task is allowed to contain. First, each
sub-task must correspond to one primitive BraneScript construct: an
import, a function definition, a conditional, a parallel block, a return
statement, and so on, all this deliberately scoped to what the \emph{workflow}
has to (and can) express, not what a package function computes internally: the
decomposition prompt explicitly forbids sub-tasks such as ``compute risk
score'' or ``validate input'', since those describe logic already
implemented inside a package rather than something the generated
BraneScript itself needs to construct from scratch. This keeps decomposition aligned
with the retrieval it feeds: a sub-task describing internal package
logic would only ever retrieve documentation about a function's
\emph{behaviour}, not the BraneScript needed to \emph{call} it. Second, each
sub-task is capped at 5--15 words and explicitly phrased as ``a BraneScript
manual search query'' rather than a natural-language description, so the
same text can be used, unmodified, as the query embedded for retrieval:
decomposition and retrieval end up sharing one representation of a sub-task.


One consequence is worth naming explicitly: decomposition never produces
code, as its rules explicitly ask for none, capping output at a bare numbered list
of tasks.

\section{Retrieval Design}\label{s:retrieval-design}

The sub-tasks produced by decomposition (Section~\ref{s:decomposition}) are
what this stage retrieves against. Gap G2 identified in Chapter~\ref{s:gaps}
requires grounding natural-language translation in the concrete capabilities of the target WMS;
how well this grounding works in practice, and whether richer context measurably improves
translation accuracy, is what \textbf{RQ2} tries to answer.
This capability surface has two very different shapes, and the design treats
them differently rather than folding both into a single retrieval mechanism:

\begin{itemize}
    \item \textbf{BraneScript syntax}, like control flows, function declarations,
        typed \texttt{class} declarations, the \texttt{parallel},
        \texttt{on} constructs and more, are fixed by the language and identical for
        every intent regardless of which packages it invokes. The full
        reference has been synthesized to be small enough to be injected as static text into
        every generation prompt rather than retrieve. Because the reference
        does not vary with the intent, a retriever over it could only return
        a subset of a fixed, already-small document: it would save at most a
        few hundred tokens per prompt, at the risk of omitting a construct
        the intent actually needs (e.g., dropping \texttt{parallel} on an
        intent that requires multi-site execution) if the retriever ranks it
        as irrelevant. This design choice was made to increased the reliability of the system, 
        as always including the full reference removed that failure mode entirely.
    \item \textbf{Package and dataset context}, by contrast, grows without
        bound as more packages are registered with Brane, and only a small
        fraction of it is relevant to any single intent. The retrieval buys a
        real reduction in prompt size here without the correctness risk
        above. The \textbf{package/dataset retriever} indexes function
        signatures, parameter types, and usage examples for each package
        registered with Brane, plus metadata about registered datasets, and
        proceeds in two stages: it first scans the intent for a known
        package or dataset name (or one of its aliases, see below) and, on a
        match, retrieves \emph{every} chunk for that package via a metadata
        filter, guaranteeing complete API coverage rather than a probabilistic
        top-$k$ selection; only when no name is detected does it fall back to
        embedding similarity search, before deduplicating the combined result. In
        this thesis, this corpus is built from the seven packages used as the
        evaluation testbed (Chapter~\ref{s:evaluation}), but the retriever
        itself is not specific to those packages.
\end{itemize}

The retriever's index is built to grow one package at a time rather than
being rebuilt whole. While treating package registration as a full rebuild would
have been the simpler alternative, it was rejected because it ties the cost of adding a single
package to the size of the entire catalogue, linearly increasing the time and compute 
required to register new packages as the system scales. 
For this reason, a newly registered package is embedded and inserted into the existing
collection on its own, leaving every previously indexed package untouched while keeping 
the system scalable and responsive to new packages being added.

\subsection{Chunking Granularity}

An earlier design indexed each package's \texttt{container.yml} specification
as a single chunk. This was found to
frequently cause the model to select the wrong function within a package when
a package exposed multiple similar functions. The design was revised to
\emph{per-function chunking}, where every function in every package's specification
becomes its own retrievable chunk.
This small change increases the precision of retrieval at the cost of a larger
index, and was adopted because function-level ambiguity was the dominant source
of retrieval-induced generation errors observed during development.

\subsection{Package Aliasing and Hybrid Retrieval}

Scientists describe their data using domain vocabulary (``heart disease data'',
``epidemic curve'') rather than the internal package names
(\texttt{healthcare}, \texttt{epidemics}) that BraneScript requires, so a
name scan against package names alone would miss most intents. For this reason, each package
directory can define a \texttt{keywords.txt} file listing such domain terms, and
the retriever loads these aliases from the collection's metadata at startup.
The name-scan stage described above therefore matches on package/dataset names \emph{or} their
aliases before falling back to similarity search. This hybrid was chosen over
pure similarity search because similarity search alone occasionally retrieved
a semantically related but wrong package (e.g., \texttt{statistics} instead of
\texttt{epidemics} for epidemic-curve intents), whereas exact name or alias
matches, when available, are unambiguous.

\section{Prompt and Feedback Design}\label{s:prompt-design}

The generation prompt combines the original intent, the retrieved
package/dataset context, the static BraneScript syntax reference, and a
system prompt encoding BraneScript-specific semantic rules not inheritable by the language syntax definition (e.g., that empty strings
are invalid arguments, that \texttt{commit\_result} must be called
exactly once per named output, and that the model shall only return Branescript valid syntax).
These rules were derived iteratively from
observed failure modes rather than specified a priori, reflecting the
practical reality that the invalid-output space of a fine-grained DSL is best
characterised empirically. Both the system prompt and the user-message
template are defined once and reused in the pipeline.

\subsection{Failure Modes and Corrective Feedback}
When a candidate workflow fails validation, the design reuses the same
prompt-construction stage for a bounded number of retries
(\texttt{MAX\_RETRIES}~$=3$): the failing script and
the error that caused it to fail are injected into the next prompt as
context, rather than starting generation from scratch. \texttt{MAX\_RETRIES}
is deliberately small and fixed rather than unbounded: if a model has not
produced a valid script after three corrective attempts, the failure is
treated as a genuine generation failure to be recorded and reported
(Chapter~\ref{s:evaluation}), rather than something an arbitrarily long retry
budget should paper over.

Compiler and runtime error messages from Brane itself are injected verbatim,
rather than replaced with a generic ``previous attempt failed'' message,
since Brane already reports which line and construct is at fault. The
quality of this feedback matters directly for retry success: a message
naming the exact undefined variable or type mismatch gives the model a much
stronger signal to correct against than a vague failure would, so the more
descriptive Brane's own error output is, the more likely the next generation
attempt is to succeed.

\section{Typed Data Representation}\label{s:typed-data}

BraneScript functions that return structured data represent it as tagged class
instances (\texttt{["ClassName", \{field: value, ...\}]}) rather than as
opaque JSON-encoded strings passed between functions. This design decision was
adopted after observing that requiring the model to construct and parse ad hoc
JSON strings inside BraneScript led to substantially more hallucinated or
malformed intermediate values; typed class instances give the model a
smaller, more constrained surface to generate correctly than free-form JSON
would.

Letting the model serialise structured values as backslash-escaped JSON strings (e.g.,
\texttt{let p := "\{\textbackslash"iterations\textbackslash": 100\}";}) 
was the specific failure pattern the fine-tuned models kept reverting to,
likely because it is the natural way to represent structured data in the
general-purpose code the base models were originally trained on. The design
responds to this on two levels: every structured value a package function
accepts or returns is expressed as a declared \texttt{class}, giving a
well-typed alternative wherever a JSON string might otherwise be used; and a
dedicated check inspects each candidate script for this specific
antipattern, triggering a retry whose corrective prompt shows the wrong
pattern directly next to the correct typed-class equivalent.

The same typed representation is also why the design treats certain Brane
values -- \texttt{Data} references to registered datasets and
\texttt{IntermediateResult} values produced by package functions -- as opaque
to BraneScript itself: the model is specifically asked to \textit{not} inspect or reconstruct
their contents from within a script, only to pass them by reference between
function calls and, when a result should be persisted, name it once via
\texttt{commit\_result}. Keeping these values opaque removes an entire class
of intermediate outputs the model would otherwise have to represent
correctly, at the cost of the model being unable to branch on their contents
inside a script.

\section{Dataset Design}\label{s:dataset-design}

The training and retrieval corpus is built around the seven packages
described in Chapter~\ref{s:implementation}. An intent/BraneScript pair is
included in the fine-tuning corpus only if it is either absent from the
recorded batch-execution results (Chapter~\ref{s:implementation}), or present
there with a successful exit code; a pair recorded as having failed
execution is dropped. Training on every authored example regardless of
outcome was considered and rejected, since it would have taught the model,
by direct example, that scripts known not to run are acceptable targets to
reproduce, precisely at the one stage, training, where such an error is
hardest to detect later. The resulting corpus is split into
training and validation sets with a fixed
$85/15$ split and a fixed random seed, so that different fine-tuning runs
compared in Chapter~\ref{s:evaluation} are trained and validated on the exact
same partition rather than one that varies run to run.

A more subtle problem surfaced specifically around the \texttt{datetime}
package: nearly every one of its functions is, by design, a thin wrapper
around the machine's current wall-clock time, so its output necessarily
differs between the moment a training example's reference output was
authored and any later moment the same BraneScript is re-executed. Early
evaluation runs surfaced this directly: a generated script that was
functionally identical to the reference would still be marked as an output
mismatch, since the reference's stored timestamp could never match a
timestamp produced at execution time, regardless of the model's actual
competence, silently deflating the output-match rate for a reason entirely
unrelated to translation quality. Since no BraneScript-level fix can make
wall-clock time reproducible, the design instead removed the affected
examples from the corpus outright, alongside stripping incidentally
non-deterministic timestamp fields returned by other packages (e.g.\ a
generation timestamp inside the epidemics package's report output), so that
the fine-tuning and output-match evaluation metrics reported in
Chapter~\ref{s:evaluation} only ever compare against genuinely reproducible
ground truth.

\subsection{Evaluation Metrics}\label{s:eval-metrics}

Three metrics are used throughout Chapter~\ref{s:evaluation} to measure
translation quality.

\emph{Compile rate} is the fraction of generated scripts that pass the
BraneScript compiler without error. BraneScript is statically typed, so any
undefined variable, undefined function call, or syntax error is surfaced at
compile time. A script that does not compile cannot be executed and is
therefore of no practical use.

\emph{Execution rate} is the fraction of generated scripts that execute to
completion without a runtime error. A script may compile successfully but fail
at runtime if it invokes a function with incorrect argument types, references
a dataset name not present in the registry, or triggers an error inside a
Brane package container.

\emph{Output match rate} is the strictest metric: the fraction of all test
examples for which the generated script's standard output is an exact string
match of the reference output. Matching is case-sensitive and order-sensitive;
outputs that contain the correct values but in a different order, or with
different formatting, are marked incorrect. The metric is computed over all
examples including those without a stored reference output, making it the most
conservative measure of end-to-end correctness.

\section{Fine-Tuning Design}\label{s:fine-tuning-design}

Several design decisions shape the fine-tuning approach used to produce
domain-adapted alternatives to the base LLMs.

\textbf{Parameter-efficient fine-tuning (QLoRA).} Full fine-tuning of a 27B
model is infeasible within the available HPC allocation. The design instead
uses 4-bit quantisation combined with Low-Rank Adaptation (LoRA), which
updates a small number of adapter parameters while keeping the base weights
frozen and quantised (exact quantisation and adapter hyperparameters are
given in Chapter~\ref{s:implementation}). This makes fine-tuning the largest
evaluated model (Qwen3.6-27B) tractable on the available GPU hardware, at the
cost of a lower-rank approximation of the full fine-tuning update.

\textbf{Full-coverage LoRA targeting.} The adapter is attached to every
linear projection in each transformer block -- the four attention
projections (query, key, value, output) and the three feed-forward
projections (gate, up, down) -- rather than to attention alone, which is a
common narrower default. BraneScript is a small, rigid DSL whose difficulty
for the base model is not primarily about attending to the right tokens in
context (what the attention projections adapt) but about reproducing exact
syntax, keyword names, and package call signatures token by token, which is
governed at least as much by the feed-forward projections; restricting
adaptation to attention only was rejected because it would have left the
feed-forward pathway -- the more likely source of the syntactic and
package-signature errors this thesis specifically aims to reduce -- frozen. The
rank is kept the same across every attention and feed-forward matrix rather
than allocated unevenly, and the LoRA scaling factor $\alpha$ is fixed at
twice the rank throughout, following the widely used convention that keeps
the effective adaptation strength comparable as rank is changed --
avoiding introducing a second hyperparameter that would otherwise need its
own justification and tuning budget alongside rank itself.

\textbf{Training-inference prompt parity.} The fine-tuning corpus is not
built from raw text but from the exact chat-formatted message list that
\texttt{prepare\_dataset.py} constructs by calling the same
\texttt{build\_user\_message} function used at inference time
(Section~\ref{s:prompt-design}): a \texttt{system} turn carrying the
BraneScript rules and syntax reference, a \texttt{user} turn carrying the
intent and retrieved package/dataset context, and an \texttt{assistant} turn
holding the reference BraneScript itself, standing in for the reply the
model is being trained to produce. The training script itself does no
prompt construction of its own; it only applies the tokenizer's chat
template to whatever messages the dataset already contains. This was a
deliberate consequence of centralising prompt construction in
\texttt{src/prompts.py} rather than a separate design decision made inside
\texttt{train.py}: since fine-tuning teaches the model to associate the
reference BraneScript with the exact prompt text surrounding it, routing
training data through the identical function used by the pipeline and the
evaluation script guarantees that whatever the model learns from is the same
distribution of prompts it is later evaluated and used on, rather than a
separately maintained approximation of it.

\textbf{Response-only loss masking.} An early version of the training loop
computed the language-modelling loss over the entire prompt, including the
system prompt and retrieved context, not just the target BraneScript. This
diluted the gradient signal for the stop token and caused generation to
continue past the intended end of the script at inference time. The design was
corrected to mask the loss so that it is computed only over the response
tokens (the reference BraneScript), which is the standard approach for
instruction-style fine-tuning and directly fixed the over-generation problem
observed during development.

\textbf{Model-agnostic training across scales.} The training script reads the
base model from an environment variable rather than hard-coding it, so the
same training procedure -- data, prompt template, LoRA configuration, loss
masking -- applies unchanged to every model size. This is what makes the
three-way comparison in Chapter~\ref{s:evaluation} (Qwen3.5-4B, Qwen3.5-9B,
and Qwen3.6-27B, each fine-tuned independently) a precise and reliable controlled comparison of
model scale.

\textbf{Resumable training and checkpoint selection.} HPC allocations are
time-bounded and jobs can be pre-empted before training completes, so the
design checkpoints periodically and, on restart, resumes automatically from
the latest checkpoint for that run rather than restarting from the base
model (an explicit flag is available to force a clean restart when needed).
Checkpoints are evaluated against the held-out validation split at the same
interval, and the checkpoint kept at the end of a run is the one with the
lowest validation loss, not necessarily the last one written -- guarding
against the case where additional epochs overfit the training split rather
than continuing to improve. Each run also writes its own configuration
(base model, LoRA settings, learning rate, dataset sizes) alongside its
checkpoints, so that a given adapter's provenance is not dependent on
external notes.

\section{Intent Provenance}\label{s:design-provenance}

In Section~\ref{s:provenance} we defined \emph{intent provenance} the mapping
from a natural-language request to the workflow it produced, alongside the
standard retrospective record of how that workflow executed. This type of provenance is identified as a gap no
existing provenance format addresses (gap G1). This thesis also aims to contribute to a design that records this
mapping for every submitted request. The pipeline,
when run with data collection enabled, writes a per-request record containing the same fields: the original intent,
the generated BraneScript, the model that produced it and which retry
attempt it was, the pass/fail verdict, the specific error type on failure
(a JSON-decoding antipattern, a Brane compilation failure, or a runtime
failure), and the captured
\texttt{stdout}/\texttt{stderr}/exit code (\texttt{TrainingCollector} in
\texttt{training\_collector.py} for the command-line tool;
\texttt{save\_to\_dashboard} in \texttt{job\_watcher.py} for the
dashboard/Snellius path). Every generation event therefore leaves behind an
explicit, machine-readable link between the intent entity, the generation
activity, and the resulting workflow entity.

These per-request records are serialised as a PROV-DM document by
\texttt{prov\_export.py}, built on the Python \texttt{prov} library. For this, we
defined a small vocabulary that specialises PROV-DM's existing entities and activities
for this system, in the same way CWLProv and ProvONE specialise them for
their own domains (Section~\ref{s:provenance}). An intent is a
\texttt{prov:Entity} \texttt{used} by a \texttt{TranslationActivity}, which
\texttt{wasGeneratedBy}-relates to a \texttt{BraneScriptWorkflow} entity and
is \texttt{wasAssociatedWith} the model that produced it; when the workflow
is executed, an \texttt{ExecutionActivity} \texttt{wasInformedBy} that
translation \texttt{used}s the workflow entity and \texttt{wasGeneratedBy}-relates
to an \texttt{ExecutionResult} entity carrying the exit code and verdict.
Entities and activities are identified by qualified names derived from the
originating request's id (e.g.\ \texttt{bt:workflow-<id>}), so that exporting
many records into one document folds shared nodes -- the same workflow
referenced from two different records -- into a single node rather than
duplicating it, instead of requiring a separate reconciliation step.
Figure~\ref{fig:intent-provenance} shows this graph for a single request,
together with the cache-hit extension described next.

\begin{figure}[htbp]
    \centering
    \resizebox{\textwidth}{!}{%
    \begin{tikzpicture}[
        entity/.style={ellipse, draw=black, thick, fill=yellow!55,
            minimum width=2.6cm, minimum height=1.1cm, align=center, font=\small},
        activity/.style={rectangle, draw=black, thick, fill=blue!25,
            minimum width=2.6cm, minimum height=1.1cm, align=center, font=\small},
        agent/.style={regular polygon, regular polygon sides=5,
            shape border rotate=90, draw=black, thick, fill=orange!55,
            minimum size=2.1cm, align=center, font=\small, inner sep=1pt},
        rel/.style={-{Latex[length=2.2mm]}, thick},
        relalt/.style={-{Latex[length=2.2mm]}, thick, dashed, teal!70!black},
    ]
        % Main generation + execution row
        \node[entity]                       (intent)   at (0,0)   {Intent\\\emph{(Entity)}};
        \node[activity]                     (transl)   at (3,0)   {Translation\\\emph{(Activity)}};
        \node[entity]                       (workflow) at (6,0)   {Workflow\\\emph{(Entity)}};
        \node[activity]                     (exec)     at (9,0)   {Execution\\\emph{(Activity)}};
        \node[entity]                       (result)   at (12,0)  {Result\\\emph{(Entity)}};
        \node[agent]                        (agent)    at (3,-2.6) {Model\\\emph{(Agent)}};

        \draw[rel] (transl) -- node[above, font=\scriptsize]{used} (intent);
        \draw[rel] (workflow) -- node[above, font=\scriptsize]{wasGeneratedBy} (transl);
        \draw[rel] (transl) -- node[left, font=\scriptsize]{wasAssociatedWith} (agent);
        \draw[rel] (workflow) to[bend left=18] node[right, font=\scriptsize, xshift=1mm]{wasAttributedTo} (agent);
        \draw[rel] (exec) -- node[above, font=\scriptsize]{used} (workflow);
        \draw[rel] (exec) to[bend right=25] node[above, font=\scriptsize]{wasInformedBy} (transl);
        \draw[rel] (result) -- node[above, font=\scriptsize]{wasGeneratedBy} (exec);

        % Cache-hit row, below
        \node[entity, fill=yellow!30]        (nintent) at (0,-5.2) {New Intent\\\emph{(Entity)}};
        \node[activity, fill=blue!12]        (lookup)  at (3,-5.2) {Cache\\Lookup\\\emph{(Activity)}};
        \node[entity, fill=yellow!30]        (served)  at (6,-5.2) {Reused\\Workflow\\\emph{(Entity)}};

        \draw[rel] (lookup) -- node[below, font=\scriptsize]{used} (nintent);
        \draw[rel] (served) -- node[below, font=\scriptsize]{wasGeneratedBy} (lookup);
        \draw[relalt] (nintent) -- node[left, font=\scriptsize]{alternateOf} (intent);
        \draw[relalt] (served) -- node[right, font=\scriptsize]{wasDerivedFrom} (workflow);
    \end{tikzpicture}}
    \caption{The PROV-DM graph produced for one request (solid, top row and
        agent) and for a cache hit against it (dashed, bottom row). The
        \texttt{Workflow} entity is typed \texttt{BraneScriptWorkflow} and
        the \texttt{Result} entity \texttt{ExecutionResult} in the exported
        document; both are shortened here for space. A cache hit does not
        run a new \texttt{TranslationActivity}: its intent is
        \texttt{alternateOf} the one that was already served, and the
        workflow it reuses \texttt{wasDerivedFrom} the original, rather than
        being generated independently.}
    \label{fig:intent-provenance}
\end{figure}

In fact, a cache hit so far leaves no trace of its own: the pipeline returnes the
cached script directly without writing a record, so the
intent-reproducibility relationship described in Section~\ref{s:provenance}
had nothing to serialise. We try to \texttt{TrainingCollector.log\_cache\_hit} close
this gap by recording, for every cache hit, the new intent, the id of the
run whose cache entry matched it, and the similarity score; the exporter then
with this implementation, on a fresh generation now also which run produced the entry it is
caching is being recorded, rather than leaving that link empty. With both sides now recorded, the
export represents a cache hit as its own small graph: the new intent is
\texttt{alternateOf} the original one, and the workflow served in response to
it \texttt{wasDerivedFrom} the workflow entity of the run that originally
generated it, rather than from a fresh translation activity, exactly the
relationship Section~\ref{s:provenance} identifies as missing from every
existing provenance format. The exporter can serialise a single request or
merge every run and cache hit under \texttt{TRAINING\_DATA\_DIR} into one
document, in PROV-N, PROV-JSON, PROV-XML, or a Graphviz rendering for
inspection. What is not being attempted in this thesis is a queryable provenance store: the
export is a file produced on demand, not a graph database a researcher can
query directly, and stable identifiers are scoped to a single run rather
than persisting across the lifetime of a deployment; both are left as future
work (Chapter~\ref{s:discussion}).

\section{Semantic Cache Design}\label{s:semantic-cache-design}

Recording intent provenance, as described above, is not by itself enough to
make the system reproducible: the generation activity it records is
performed by a non-deterministic LLM, so submitting the same intent again --
or the same request paraphrased differently -- could still produce a
different workflow entity each time, with no relation recorded between the
two events. The semantic cache is the mechanism that closes this gap (the
intent-reproducibility gap, G3, Section~\ref{s:provenance}): the pipeline
queries the cache before invoking the LLM at all, and if the same or a
paraphrased intent has been served before, the cached BraneScript is
returned deterministically instead of triggering a fresh, potentially
divergent generation event. Only on a cache miss does the pipeline fall back
to LLM generation. This is a single fixed lookup-then-generate ordering, not
an adaptive or learning mechanism: the cache does not become more or less
trusted over time, and no component of the pipeline tracks confidence across
requests; a request that can be served from the cache simply never reaches
the LLM, which also means it is never re-generated at a fraction of the
cost. The cache is consulted by both the command-line tool and the
dashboard, but is deliberately excluded from the offline evaluation pipeline
used to produce the results in Chapter~\ref{s:evaluation}: a cache hit there
would let a previously generated correct answer masquerade as freshly
generated output, inflating the reported compile and execution rates without
reflecting the quality of the model actually being evaluated.

An exact-string-match cache -- looking up the intent verbatim rather than by
embedding similarity -- was considered and rejected as the primary mechanism.
It avoids the false-positive risk discussed below, but only ever hits on the
literal same sentence being submitted twice, defeating the purpose in a
system whose stated premise (Chapter~\ref{s:overview}) is that a researcher
can phrase a request however feels natural to them; two researchers
describing the same analysis in different words would never share a cache
entry under exact matching.

This is a genuine trade-off, not a settled solution: the same embedding
space that enables useful generalisation across paraphrases also maps intents
that differ only in a specific parameter (e.g., a numeric threshold or a
dataset name) to very high similarity, since sentence embeddings capture
propositional meaning rather than exact parameter values. The cache threshold
is therefore an explicit precision--recall operating point that must be
selected empirically: raising it reduces false positives at the cost of hit
rate, while lowering it increases coverage at the cost of precision. The
empirical characterisation in Chapter~\ref{s:evaluation} identifies 0.92 as
a reasonable default, but the trade-off itself cannot be eliminated through
threshold tuning alone.

Characterising that operating point requires intents that are genuine
paraphrases of an existing cache entry, as they should just be worded differently while still referring
to the same underlying request. The initial curated dataset does not provide
on its own, since it was meant to be varied and sparse rather than to contain
near-duplicates intents. \texttt{scripts/generate\_paraphrases.py} closes this gap by
prompting an LLM (Qwen3.5-4B) to reword every training intent several times
while keeping its reference BraneScript unchanged, writing the result to
\texttt{data/training/paraphrases.jsonl}. Each paraphrase is only ever
compared against the cache, never merged into the fine-tuning corpus itself
(Section~\ref{s:dataset-design}): its purpose is exclusively to sweep the
similarity threshold and measure, for each candidate value, how often a
paraphrase is matched to any cache entry at all and how often it is matched
to the \emph{correct} one, giving the precision--recall figures reported in
Section~\ref{s:eval-cache} from which the deployed threshold was chosen.

A cache entry is written only once a candidate script has cleared syntax and
semantic checking (or, when execution is requested, only once it has also
executed successfully), so a script that the pipeline itself rejected can
never be served back as a cache hit later. Because the threshold is a
similarity cutoff rather than an exact match, the design also has to account
for the case where a cached entry turns out to be wrong for reasons outside
the pipeline's own checks (e.g., an underlying package or dataset changes
after the script was cached, so a previously correct answer no longer holds).
Rather than leave stale entries to accumulate indefinitely, the cache exposes
an explicit \texttt{invalidate} operation, reachable through the dashboard
backend's \texttt{/api/cache/invalidate} endpoint, that removes every entry
within the similarity threshold of a given intent, giving a way to correct
the cache without needing direct database access. Each entry additionally tracks a hit
counter, incremented on every cache hit, giving a simple, low-cost measure of
how much reuse a given cached answer is actually getting.

\subsection{Cache Evaluation Metrics}\label{s:cache-metrics}

Three metrics are used to characterise the cache threshold in
Section~\ref{s:eval-cache}.

\emph{Hit rate} is the fraction of queries for which the cache returned any
result, i.e.\ a cached entry whose similarity to the query exceeded the
threshold.

\emph{Correct rate} is the fraction of all queries for which the returned
BraneScript matched the reference for that intent. A query that misses the
cache counts as incorrect for the purpose of this metric.

\emph{False-positive rate} is the fraction of queries where the cache
returned a result but it did not match the reference, indicating that a
semantically distinct intent was incorrectly retrieved above the threshold.
